%%%%%%%%%%%%%%%%%%%%%%%%%%%%%%%%%
%
%       EXERCÍCIO
%
%%%%%%%%%%%%%%%%%%%%%%%%%%%%%%%%%

\ifdefstring{\atividade}{prova}{%
    \ifdefstring{\modo}{objetivo}{%
        \renewcommand{\valorquestao}{\ValorQObj\ ponto}
    }{%
        \renewcommand{\valorquestao}{\ValorQDisc\ pontos}
    }%
}%

\begin{exercicioBanco}[\valorquestao]
Três cidades \(A\), \(B\) e \(C\) situam-se ao longo de uma estrada reta; \(B\) situa-se entre \(A\) e \(C\) e a distância de \(B\) a \(C\) é igual à metade da distância de \(A\) a \(B\). Um encontro foi marcado por 3 moradores, um de cada cidade, em um ponto \(P\) da estrada, localizado entre as cidades \(B\) e \(C\) e à distância de 180 km de \(A\). Sabendo-se que \(P\) está 30 km mais próximo de \(C\) do que de \(B\), determinar a distância que o morador de \(B\) deverá percorrer até o ponto de encontro.

% Define as alternativas
\newcommand{\alternativas}{%
    \begin{center}
        \begin{tabularx}{\textwidth}{XXXXX}
            (a) \(30\) km. &
            (b) \(40\) km. &
            (c) \(50\) km. &
            (d) \(60\) km. &
            (e) \(70\) km.
        \end{tabularx}
    \end{center}
}

% Define a resposta correta
\newcommand{\resposta}{D}

% Lógica condicional para exibição
\ifdefstring{\atividade}{lista}{%
    \alternativas
    \vspace{0.5em}
    
    \noindent\textbf{Resposta:} letra \textbf{\resposta}.
}{%
    \ifdefstring{\modo}{objetivo}{%
        \alternativas
    }{%
        % modo = discursiva → não mostra alternativas
    }
}
\end{exercicioBanco}

